%% bab5.tex
%% BAB 5 - KESIMPULAN DAN SARAN
%% Format sesuai Panduan Penyusunan Tesis ITS 2021

\chapter{KESIMPULAN DAN SARAN}

\section{Kesimpulan}

Pada sub-bab ini dituliskan kesimpulan hasil penelitian atau kesimpulan tesis. Kesimpulan harus ditulis berdasarkan hasil penelitian, pembahasan, dan temuan yang telah ditulis pada bab sebelumnya yang tentu saja disesuaikan dengan tujuan penelitian atau tesis. Jangan menyimpulkan sesuatu yang tidak ada di dalam pembahasan yang telah dibuat.

Berdasarkan hasil penelitian dan pembahasan yang telah dilakukan, dapat ditarik kesimpulan sebagai berikut:

\begin{enumerate}
    \item Kesimpulan pertama sesuai dengan tujuan penelitian pertama. Jelaskan temuan utama yang menjawab rumusan masalah pertama.
    
    \item Kesimpulan kedua sesuai dengan tujuan penelitian kedua. Jelaskan temuan utama yang menjawab rumusan masalah kedua.
    
    \item Kesimpulan ketiga sesuai dengan tujuan penelitian ketiga. Jelaskan temuan utama yang menjawab rumusan masalah ketiga.
\end{enumerate}

\section{Saran}

Pada sub-bab ini dituliskan saran yang diusulkan oleh penulis. Saran yang dibuat harus berdasarkan pembahasan serta kesimpulan yang telah dibuat. Jangan menyarankan sesuatu yang berada di luar jangkauan pembahasan dan kesimpulan yang dibuat.